\SourceLeaf{112}
die Vorschriften zur vollständigen Ausgleichung der Beobachtungen immer einerlei Resultat geben, sie mögen auf die ursprünglichen Beobachtungen selbst, oder auf die bereits einstweilen \textit{unvollständig} ausgeglichenen Beobachtungen angewandt werden, in so fern dieser Begriff in der in der Abhandlung näher bestimmten Bedeutung genommen wird, unter welcher, als specieller Fall, derjenige begriffen ist, wo mit den Beobachtungen schon eine zwar vorschriftsmässig ausgeführte, aber nur einen Theil der Bedingungsgleichungen berücksichtigende Ausgleichung vorgenommen war.

Den letzten Theil der Abhandlung machen ein paar mit Sorgfalt ausgearbeitete Beispiele der Anwendung der Methode aus, die theils von den geodätischen Messungen des Generals von \textsc{Krayenhoff}, theils von der vom Verfasser selbst im Königreich Hannover ausgeführten Triangulirung entlehnt sind, und die dazu dienen können, sowohl die Anwendung dieser Theorie mehr zu erläutern, als auch manche, dergleichen Messungen betreffende, Umstände überhaupt in ein helleres Licht zu stellen.

Die trigonometrischen Messungen gehören ganz besonders in das Feld, wo die Wahrscheinlichkeitsrechnung Anwendung findet, und namentlich in derjenigen Form Anwendung findet, die in der gegenwärtigen Abhandlung entwickelt ist. Gerade hier ist es Regel, dass mehr beobachtet wird, als unumgänglich nöthig ist, und dass so die Messungen einander vielfältig controlliren. Nur durch die Benutzung der strengen Grundsätze der Wahrscheinlichkeitsrechnung kann man von diesem Umstande den Vortheil ganz ziehen, der sich davon ziehen lässt, und den Resultaten die grösste Genauigkeit geben, deren sie fähig sind. Ausserdem aber geben jene Grundsätze zugleich das Mittel, die Genauigkeit der Messungen selbst, und die Zulässigkeit der darauf gegründeten Resultate zu bestimmen. Endlich dienen sie dazu, bei der Anordnung des Dreieckssystems, aus mehreren, unter denen man vielleicht die Wahl hat, das zweckmässigste auszuwählen. Und alles dieses nach festen sichern Regeln, mit Ausschliessung aller Willkürlichkeiten. Allein sowohl die sichere Würdigung, als die vollkommenste Benutzung der Messungen ist nur dann möglich, wenn sie in reiner Authenticität und Vollständigkeit vorliegen, und es wäre daher sehr zu wünschen, dass alle grösseren auf besondere Genauigkeit Anspruch machenden Messungen dieser Art immer mit aller nöthigen Ausführlichkeit bekannt gemacht werden möchten. Nur zu gewöhnlich ist das Gegentheil, wo nur Endresultate für die einzelnen gemes\SourceLeaf{113}senen Winkel mitgetheilt werden. Wenn solche Endresultate nach richtigen Grundsätzen gebildet werden, indem man durchaus alle einzelnen Beobachtungsreihen, die nicht einen durchaus unstatthaften Fehler gewiss enthalten, dazu concurriren lässt, so ist der Nachtheil freilich lange nicht so gross, als wenn man etwa nur diejenigen Reihen beibehält, die am besten zu den nahe liegenden Prüfungsmitteln passen, welche die Summen der Winkel jedes Dreiecks und die Summen der Horizontalwinkel um jeden Punkt herum darbieten. Wo dies durchaus verwerfliche Verfahren angewandt ist, sei es aus Unbekanntschaft mit den wahren Grundsätzen einer richtigen Theorie. oder aus dem geheimen Wunsche, den Messungen das Ansehen grösserer Genauigkeit zu geben, geht der Maasstab zu einer gerechten Würdigung der Beobachtungen und der aus ihnen abzuleitenden Resultate verloren; die gewöhnliche Prüfung nach den Winkelsummen in den einzelnen Dreiecken, und bei den Punkten. wo die gemessenen Winkel den ganzen Horizont umfassen. scheint dann eine Genauigkeit der Messungen zu beweisen, von der sie vielleicht sehr weit entfernt sind, und wenn andere Prüfungsmittel, durch die Seitenverhältnisse in geschlossenen Polygonen oder durch Diagonalrichtungen, vorhanden sind, werden diese die Gewissheit des Daseins von viel grössern Fehlern verrathen. Umgekehrt aber. wenn die zuletzt erwähnte Voraussetzung Statt findet, und das Ausgleichen der Beobachtungen in Beziehung auf die Prüfungsmittel ohne die sichern Vorschriften der Wahrscheinlichkeitsrechnung versucht ist, wo es immer ein Herumtappen im Dunkeln bleiben muss, und grössere, oft viel grössere, Correctionen herbeiführt, als nöthig sind, kann leicht dadurch ein \textit{zu} ungünstiges Urtheil über die Messungen veranlasst werden. Diese Bemerkungen zeigen die Wichtigkeit sowohl einer hinlänglich ausführlichen Bekanntmachung, als einer auf strenge Principien gegründeten mathematischen Combination der geodätischen Messungen: sie gelten aber offenbar mehr oder weniger bei Beobachtungen jeder Art, astronomischen, physikalischen u.\,s.\,w., die sich auf das Quantitative beziehen, insofern die Mannigfaltigkeit der dabei Statt findenden Umstände zu wechselseitigen Controllen Mittel darbietet.
\par\medskip\centerline{\rule{.3\linewidth}{.3pt}}
